% Faithful PDF text layer: cmap (before fontenc) attaches real ToUnicode
% maps so ligatures (fi, fl, ff...) copy/extract as proper letters —
% "fortifications", not "forti?cations". lmodern guarantees Type 1 fonts.
% Display is unaffected.
\usepackage{cmap}
\input{glyphtounicode}
\pdfgentounicode=1

\usepackage[utf8]{inputenc}
\usepackage[LGR,T1]{fontenc}
\usepackage{lmodern}
\usepackage[main=english,hebrew]{babel}
\usepackage[margin=2.5cm, headheight=14pt]{geometry}
\usepackage{hyperref}
\usepackage{graphicx}
\usepackage{amsmath}
\usepackage{amssymb}
\usepackage{booktabs}
\usepackage{tabularx}
\usepackage{fancyhdr}
\usepackage{float}
\usepackage{tikz}
\usetikzlibrary{positioning, arrows.meta}

\hypersetup{
    colorlinks=true,
    linkcolor=blue,
    urlcolor=blue,
    citecolor=blue
}

% Headers/footers are forced into the English language context: when a page
% break lands at a Hebrew RTL block, the next page ships while babel is still
% in Hebrew mode, so an unprotected header prints mirrored and with the
% Hebrew font's digit glyphs (which look Arabic-Indic). \foreignlanguage
% resets both direction and font; safe here because headers are English-only.
\pagestyle{fancy}
\fancyhf{}
\fancyhead[L]{\foreignlanguage{english}{\small\textit{\BookTitle}}}
\fancyhead[R]{\foreignlanguage{english}{\small\textit{Intelligent Agents --- Ex03}}}
\fancyfoot[C]{\foreignlanguage{english}{\thepage}}
\renewcommand{\headrulewidth}{0.4pt}
\renewcommand{\footrulewidth}{0.4pt}

\fancypagestyle{plain}{
  \fancyhf{}
  \fancyfoot[C]{\foreignlanguage{english}{\thepage}}
  \renewcommand{\headrulewidth}{0pt}
}
